\section{Open questions}

The following five questions remain genuinely open after this replication. Each is anchored in a specific claim, gap, or systematic that the current work either identified but did not resolve, or explicitly deferred as out of scope. Each carries a concrete, executable next step (not ``more analysis'').

\begin{enumerate}

\item \textbf{Is the WTe$_2$ 32-qubit downfolded result bond-dimension-limited?}
\emph{Basis.} The current replication reproduces the Ca$_2$CuO$_3$ 20-qubit result at $\sim$0.1~meV but skips WTe$_2$ (Fock dim $\sim 10^9$, beyond direct ED). The paper's WTe$_2$ VQE uses truncated MPS bond dimension $\chi=512$ without a $\chi$-convergence sweep in the main text; the reported VQE $\Delta_{\text{exc}} = 0.379$ vs.\ DMRG $\Delta_{\text{exc}}=0.640$ (a nearly 2$\times$ discrepancy) may be either an ansatz-fidelity failure or a bond-dimension artifact.
\emph{Next steps.} Build an ITensor DMRG driver for the 4-band 2$\times$2 extended-Hubbard Hamiltonian using Appendix~C.2 parameters. Sweep $\chi \in \{64, 128, 256, 512, 1024\}$ and check $E_0$ convergence. Compare against a smaller 1$\times$2 or 2$\times$1 sub-lattice where quasi-exact ED gives ground truth. Report $\Delta_{\text{exc}}(\chi)$ to attribute the VQE-vs-DMRG gap.

\item \textbf{How sensitive are the reported downfolded matrices to the cRPA / active-space choice?}
\emph{Basis.} The paper picks one cRPA window and one active-space definition per material and reports one set of $t, U, V$ matrices. Downfolded-Hamiltonian methods (DUCC, cRPA, wan2respack) are known to inherit systematics from these choices, but the paper does not report a sensitivity table. The internal ED-vs-DMRG agreement here validates the compressed model, not the compression itself.
\emph{Next steps.} Re-run Wannier90 + RESPACK for Ca$_2$CuO$_3$ with three cRPA windows (paper's, $\pm$0.5~eV wider, $\pm$0.5~eV narrower). Recompute $t, U, V$ and diagonalise (ED is affordable for Ca$_2$CuO$_3$); quote $E_0$ and nearest-neighbour spin correlation for each. This bounds the downfolding truncation-sensitivity envelope on the paper's headline number without needing the quantum-hardware stack.

\item \textbf{Is the downfolding compression itself lossless, in the sense the paper claims?}
\emph{Basis.} The paper argues downfolding preserves low-energy physics but does not quote a direct comparison against a non-downfolded reference on the same material. The 0.1~meV ED-vs-DMRG agreement here validates the compressed model internally; it does not prove that the compressed model matches an ab-initio reference without an active-space truncation.
\emph{Next steps.} Use a public AFQMC package (\texttt{pyqmc} or \texttt{ipie}) on Ca$_2$CuO$_3$ with a plane-wave or Gaussian basis at the paper's DFT-PBE geometry to obtain a total energy per Cu site and a spin structure factor. Compare against the downfolded ED result rescaled per Cu site. A gap of more than $\sim$tens of meV per site would flag compression loss. Extend to WTe$_2$ excitonic order if resources permit.

\item \textbf{What ground-state accuracy is achievable on today's superconducting hardware under a realistic error-mitigation stack?}
\emph{Basis.} Table~II reports 0.999-per-gate ideal-model circuit fidelity 74.8\% but the paper is entirely classical MPS-VQE simulation; no real device closes the loop between the resource estimate and a delivered energy. Real superconducting hardware has non-Markovian noise, crosstalk, and calibration drift that a $0.999^{n_{2q,G}}$ model does not capture.
\emph{Next steps.} Compile the 20-qubit / 290-two-qubit-gate NP ansatz to an IBM Eagle-127 or Heron-156 device via Qiskit. Submit at three zero-noise scale factors, apply linear + exponential ZNE, and report the extrapolated energy vs.\ paper's 6.028~eV VQE and 6.005~eV DMRG. Repeat with 1000 M3 readout-mitigation shots per Pauli term. Quantify how many of the 23~meV VQE-vs-DMRG gap survives mitigation.

\item \textbf{Does the workflow scale into genuinely strongly-correlated regimes, or does the required MPS bond dimension blow up?}
\emph{Basis.} Ca$_2$CuO$_3$ is quasi-1D and spin-gapped; WTe$_2$ is an excitonic insulator; SrVO$_3$ is a moderately-correlated metal but the paper's downfolded model still converges. None of the three stress-tests the entanglement-scaling regime where 2D critical systems require $\chi \sim \exp(L)$ (underdoped cuprates, moir\'e heterostructures, Kondo lattices).
\emph{Next steps.} Apply the full Wannier + cRPA + MPS-VQE workflow to (i) a doped Hubbard model on a 4$\times$4 square lattice at hole densities 0.05, 0.10, 0.15 (approaching the stripe-critical regime) and (ii) the magic-angle twisted-bilayer-graphene continuum-model reduction. Track $\chi$ needed for 1~meV per site convergence vs.\ system width, benchmark against published DMRG (White--Scalapino, tBG DMRG). Monotonic $\chi$ blow-up identifies a hard failure mode; graceful growth strengthens the generality claim. Cross-benchmark against DMRG-only on the same downfolded $H$ to isolate the ansatz limitation from the downfolding limitation.

\end{enumerate}
